\chapter*{历史注记}
\phantomsection
\addcontentsline{toc}{chapter}{历史注记}

（注意——罗马数字指本注记之末所附的参考文献。）

以现代面貌出现的二次型理论，其历史很难追溯到十八 \(^{e}\) 世纪后半叶之前，而且，如我们将要看到的，它的发展首先是为了满足数论、分析与力学的需要。但是这一理论的基本概念实际上从“欧几里得”几何的开端起就已出现，并且构成其骨架。因此，若不至少概略地谈谈“初等几何”自古以来的发展，就无法追溯它的历史。当然，我们只能关注若干一般观念的演变，读者也不应指望在这里找到关于任何特定定理的历史的精确资料，对此我们只需引证专门的历史著作或教学著作（*）即已足够。不言而喻，当我们在下文谈到同一条定理在各种代数语言或几何语言中的各种可能解释时，绝不是说这些“翻译”在各个时代都像今天这样为人们所熟悉；恰恰相反，本注记的主要目的正是要表明，数学家们是如何极其渐进地意识到这些往往面貌迥异的问题之间的亲缘关系的；我们还要表明，他们在此过程中如何被引导着为古人留下的大量几何定理引入某种连贯性，并最终试图精确地划定“几何”一词应当理解成什么。

除了巴比伦人发现的求解二次方程的公式（(I)，pp. 183-189）之外，二次型理论主要概念的产生都必须在其几何伪装之下来考察。这些概念首先以距离的平方（在平面上或在三维空间中）的面貌出现，而相应的“正交性”概念则是借助于直角引入的，欧几里得把直角定义为平角的一半（《几何原本》第 I 卷，定义 10）；距离与直角这两个概念由勾股定理——欧几里得大厦的拱顶石（*）——联系在一起。角的概念在希腊数学中出现得似乎很早（希腊人很可能得之于巴比伦人，后者凭借长期的天文经验精于角的使用）。众所周知，在古典时期，只定义小于两直角的角（欧几里得的“定义”与他关于直线或平面给出的定义一样含糊而不可用）；定向的概念并未被提出，尽管欧几里得使用（既无公理又无定义）了一条直线把平面分成两个区域这一事实，并在必要处仔细地加以区分（**）。在这个阶段，平面旋转群的思想只是以很不完善的方式，通过（也是欧几里得不加解释地引入的）半直线非定向角的加法显现出来，而这种加法原则上也只在和至多等于两直角时才有定义（***）。至于三角学，它受到几何学家的轻视，而被留给测量员与天文学家；正是后者（阿里斯塔克斯、希帕恰斯，尤其是托勒密（V））建立了（平面的或球面的）直角三角形的边与角之间的基本关系，并编制了最早的数表（这些数表给出由角 \(\theta < \pi\) 在半径为 \(r\) 的圆上截得的弧的弦，换言之即数 \(2r \sin \frac{\theta}{2}\) ；更为便于使用的正弦的引入归功于中世纪的印度数学家）；在计算这些数表时，当时尚不知道的弧的加法公式被关于内接于圆的四边形的托勒密定理（可能上溯到希帕恰斯）的等价用法所代替（参见 Esp. vect. top.，第 V 章，§ 1，习题 5）。还应注意，欧几里得与海伦给出了与公式

\[
a ^ {2} = b ^ {2} + c ^ {2} - 2 b c \cos A
\]

等价的命题，它联系任意平面三角形的边与角；但由于缺乏直到十九 \(^{e}\) 世纪才出现的向量分析的思想，人们很难把这看作与度量型相伴的双线性型概念的首次露面。

位移（或运动，这两个概念的区分在古代——甚至在很久以后——并不清楚）已为欧几里得所知；但由于我们不知道的原因，他对使用它们似乎表现出明显的迟疑（例如在“三角形相等的情形”中，人们得到的印象是：他之所以使用位移的概念，只是因为未能表述出一条恰当的公理（(III bis)，t. I，pp. 225-227 及 249））；然而，正是借助于位移（绕轴旋转）的概念，他给出了旋转锥面与球面的定义（《几何原本》第 XI 卷，定义 14 与 18），阿基米德在定义旋转二次曲面时也是如此。但是，作用于整个空间的变换这一一般思想，在十八 \(^{e}\) 世纪末之前几乎不为数学思维所知（*）；

而在十七 \(^{e}\) 世纪之前，同样没有运动合成概念的踪迹，更不用说位移的合成了。这当然不意味着希腊人对图形的“正则性”与“对称性”不特别敏感——我们今天把这些性质与位移群的概念联系在一起；他们的正多边形理论、尤其是正多面体理论——他们全部数学中最出色的篇章之一——恰恰证明了相反的情形（*）。

最后，希腊数学在我们所关心的领域中的最后一项本质贡献是圆锥曲线的理论（至于二次曲面，希腊人只认识某些旋转二次曲面，而且除球面之外并未把研究推得很远）。这里值得注意的是：虽然希腊人从未有过解析几何基本原理的思想（主要是因为缺乏一种便于使用的代数），但他们在研究特殊“图形”时通常使用平面内相对于两条（甚至多于两条）轴的“纵标”（这些轴与图形密切相关，这正是他们的方法与费马和笛卡儿的方法根本不同之点之一，后者的轴是独立于所考虑的图形而固定的）。特别地，在与倍立方问题的联系中出现的最早一批（异于圆的）圆锥曲线，是由方程 \(y^{2} = ax\) 、 \(y = bx^{2}\) 、 \(xy = c\) 给出的曲线（门奈赫莫斯，欧多克索斯的学生， \(1v^{e}\) 世纪中叶）（**）；而在研究与这些曲线有关的问题时，最常用的是圆锥曲线的方程（通常相对于由一条直径与该直径和曲线的一个交点处的切线所成的两条斜轴）（而“焦点”性质只起很不显著的作用，与仅仅上溯到 \(x1x^{e}\) 世纪的学校传统可能使人相信的相反）。在这门庞大的理论中，我们在此尤须记住的是共轭直径的概念（阿基米德已经知道它），以及如今用作一点的极线之定义的那条性质，它由阿波罗尼奥斯（IV）在点位于圆锥曲线之外时给出（对他而言，极线就是从该点引出的各切线的切点所连成的直线）；从我们的观点看，这是关于异于度量型的二次型的“正交性”的两个例子，但这些概念与垂线的经典概念之间的联系在当时当然是不可能被想到的。

在笛卡儿与费马之前，几乎没有别的进展可报道；但从解析几何的一开始，二次型的代数理论就开始从它的几何外壳中显露出来：费马知道平面上一个二次方程表示一条圆锥曲线（(VII a)，pp. 100-102），并对二次曲面勾画了类似的想法（VII b）。随着十八 \(^{e}\) 世纪二维与三维解析几何的发展，（特别是联系圆锥曲线与二次曲面）出现了这一理论的两个中心问题：把二次型化归为平方和，以及寻求它关于度量型的“轴”。对圆锥曲线而言，这两个问题太初等，不会引起重要的代数进展；对任意多个变量的情形，第一个问题由拉格朗日在 1759 年联系多元函数的极大值问题解决（IX a）。但这个问题几乎立即被寻求轴的问题所掩盖，甚至在秩的不变性被表述之前就是如此（*）；至于惯性定律，它直到 1850 年前后才由雅可比（XIX b）与西尔维斯特（XX）发现，前者用与今天所用相同的推理证明了它，后者则只是把它当作几乎显然的事实加以陈述（**）。

把二次曲面化归到其轴的问题已呈现出比圆锥曲线的相应问题大得多的代数困难；最早着手研究它的欧拉不能证明特征值的实性，他在一段缺乏证明价值的概述之后承认了它（(VIII a)，pp. 379-392）（**）。即使这一点在 1800 年前后得到正确确立，仍须等到柯西才对任意变量个数 \(n\) 的型证明了相应的定理（XIV b）。也是柯西，在大约同一时间，证明了给出特征值的特征方程在直角轴的任何改变下不变（(XIV a)，p. 252）（****）；但对 \(n = 2\) 或 \(n = 3\) ，由于借助相应的圆锥曲线或二次曲面的轴对特征值所作的几何解释，这种不变性在直观上是“显然”的。此外，在关于这个课题的研究过程中，特征值的初等对称函数也自然而然地出现（带有种种几何解释，特别涉及阿波罗尼奥斯关于共轭直径的定理），特别是判别式，它（就 n = 2 而言，早在与二次方程理论的联系中便已为人所知）在 n = 3 时由欧拉首次给出（(VIII a) p. 382）；欧拉是在对二次曲面分类时遇到它的（在表达一个二次曲面没有任何无穷远点的条件时），而未提到它在直角轴改变下的不变性。但稍后，随着整系数二次型算术理论的开端，拉格朗日（就 n = 2）注意到判别式在线性但非正交的变量替换下的不变性的一个特例（(IX b)，p. 699），而高斯对 n = 3 建立了判别式在任何线性变换下的“协变性”（(XI a)，pp. 301-302）（*）。在柯西与比内证明了行列式乘法的一般公式之后，把高斯的公式推广到任意多个变量的情形便是直接的了；正是它在大约 1845 年给出了不变量一般理论的最初推动。

在希腊人那里取代位移理论的两个概念——即运动的概念与图形的“对称”的概念——之外，十七 \(^{e}\) 与十八 \(^{e}\) 世纪又加上了第三个，即直角轴变换的问题，它与这个理论实质上等价。欧拉有好几部著作致力于这个问题，他尤其力求为轴变换公式得到便于使用的参数表示。我们知道力学后来对他为此目的在 \(n = 3\) 时引入的三个角作了怎样的使用（(VIII a)，pp. 371-378）。但他并未止步于此；他在 1770 年考虑了对任意 \(n\) 的正交变换的一般问题，注意到通过引入 \(n(n - 1)/2\) 个角作为参数便可达到目的，最后，对 \(n = 3\) 与 \(n = 4\) 给出了旋转的有理表示（分别用 4 个齐次参数与 8 个由一个关系式联系起来的齐次参数表示），它们正是后来借助四元数理论得到的那些表示（参见 § 9，习题 15 与 16），而他并未指出其来源（VIII c）（**）。

另一方面，欧拉还指出了如何解析地寻求平面图形的“对称”，正是在这一场合，他实质上证明了平面位移是一个旋转，或一个平移，或一个平移加一个对称（(VIII a)，pp. 197-199）。当时力学的兴起还导致对位移的一般研究；但起初只有与连续运动相切的“无穷小”位移才是讨论的对象：在托里拆利、罗贝瓦尔与笛卡儿关于平面运动的合成与转动瞬时中心的研究中，出现的显然只有这种位移（参见第 IV 卷的历史注记，第 I-II-III 章）。转动瞬时中心由约翰·伯努利以一般的方式定义；达朗贝尔在 1749 年、欧拉在次年推广了这一概念，证明了使一点保持不动的运动存在转动的瞬时轴。有限位移的类似定理直到 1775 年才由欧拉陈述（VIII d），在这篇论文中他还同时发现旋转的行列式等于 1；次年，他证明了平面相似变换存在不动点（VIII e）。但必须等到夏尔自 1830 年起的工作（XV a），才最终得到关于有限位移与无穷小位移的一个连贯的理论。
% semantic-fragments: legacy-input-seam
\relax{}
\begin{center}
\(\ast\qquad\ast\)
\end{center}

我们于是达到了可谓几何学黄金时代的时期，它大致介于蒙日《画法几何》（1795）（X）与 F. 克莱因的“埃尔朗根纲领”（1872）（XXV b）两书的出版年代之间。我们得自几何学这一突然复兴的主要进展如下：

\begin{enumerate}
\def\labelenumi{\Alph{enumi})}
\item
  无穷远元素（点、直线或平面）的概念由德萨格于十七世纪引入（VI），但在十八世纪，它几乎只是作为对术语的滥用而出现；彭赛列恢复了这一概念的地位并系统地加以使用（XII），从而把射影空间确立为一切几何现象的一般框架。
\item
  与此同时，在蒙日、尤其是彭赛列那里，完成了向复射影几何的过渡。十八世纪中偶尔使用的虚点概念，在这里（与无穷远点概念一起）得到利用，以得出独立于实仿射几何各种“图形情形”的命题。诚然，起初为支持这些革新而提出的论证仍然十分笨拙（几何学“综合”学派的拥护者们尤其如此，在那里，坐标的使用曾被视作一种污点），但人们不能不从中辨认出——在蒙日名下的“偶然关系原理”、在彭赛列名下的“连续性原理”之下——现代代数几何中“特殊化”思想的最初萌芽 (*)。
\end{enumerate}

由这些观念所得的首批成果之一，是注意到在复射影空间中，一切非退化二次曲线（相应地，二次曲面）都具有同一本性；这引导彭赛列发现了“迷向”元素：“平面上任意放置的各个圆，”他说，“因此并非如人们乍看之下可能相信的那样彼此完全独立；它们在理想意义下于无穷远处有两个公共的虚点”（（XII），第 48 页）。再往后，他又同样引入了“脐”，即所有球面公共的无穷远虚二次曲线（（XII），第 370 页）；尽管他没有特别论及球面的迷向母线，他至少明确强调了所有二次曲面都具有实的或虚的直母线这一事实（同上书，第 371 页）(*)；他的后继者们（特别是普吕克与夏尔）对这些概念的利用甚至比他本人更多，尤其是在二次曲线与二次曲面的“焦点”性质的研究中。

\begin{enumerate}
\def\labelenumi{\Alph{enumi})}
\setcounter{enumi}{2}
\item
  点变换与变换的复合这两个概念，也同样以一般的方式得到表述，并被系统地引入作为证明手段。在位移与投影之外，当时已知的变换只有寥寥几种特殊情形：拉伊尔与牛顿使用的 \(x' = a/x\)、\(y' = y/x\) 型的某些平面射影变换，来自克莱罗与欧拉的“仿射” \(x' = ax\)、\(y' = by\)，最后是牛顿、麦克劳林与布雷肯里奇那里的某些特殊的二次变换。蒙日在其《画法几何》中展示了平面投影在三维几何中所能发挥的全部作用。在彭赛列那里，一种被使用到过分程度的系统证明方法，在于借助投影把二次曲线的性质化为圆的性质（德萨格与帕斯卡已经偶尔用过这一方法）；而为了类似地从二次曲面过渡到球面，他发明了空间中射影变换的第一个例子，即“透射”（（XII），第 357 页）；最后，也正是他引入了曲线到自身的双有理变换的第一批例子。1827 年，默比乌斯（（XIII a），第 217 页）（以及 1830 年的夏尔，独立地（XV b））定义了最一般的射影线性变换；同一时期还出现了反演（参见 §10，习题 13）以及其他类型的二次变换，对它们的研究将开创双有理变换的理论，这一理论将在十九世纪下半叶得到发展。
\item
  对偶的概念充分显露出来，并被自觉地与双线性型理论联系起来。关于二次曲线的极点与极线的理论——自阿波罗尼奥斯以来仅在德萨格与拉伊尔的著作中取得过进展——被蒙日推广到二次曲面，蒙日与其学生们一起察觉到，借此手段可以把已知的定理变换为新的结果 (*)。但是，把上述种种观察确立为一般方法——即其“配极变换”理论——并把该方法变成一种特别有效的发现工具，这一功劳仍应归于彭赛列。稍后，特别是在热尔岗、普吕克、默比乌斯与夏尔那里，一般的对偶概念从它与二次型的联系（这种联系在彭赛列那里仍过于狭窄）中脱离出来。特别是，默比乌斯在考察三维空间中对偶（由一个双线性型定义）的各种可能性时，于 1833 年发现了关于交错双线性型的对偶 (XIII b) (**)，它在十九世纪尤其以“线性丛”理论的形式（参见 §10，习题 16）得到研究，并与凯莱、格拉斯曼和普吕克约在 1860 年前后引入的“直线几何”及“普吕克坐标”相联系而发展起来。
\item
  从射影几何的最初时期起，对经典几何诸性质及其与射影空间关系的深入研究，便很快导致把这些性质划分为“射影性质”与“度量性质”；把这一区分视为当时未来将成为现代结构概念的东西的最清晰表现之一，大概并不算夸张。但首先引入这一区分及这套术语的彭赛列，已经意识到联系这两类性质的东西；他在其《论著》中处理关于角的诸问题时说，角的性质“似乎不属于我们称之为射影的那类性质\ldots，然而它们却以如此简单的方式”，他说，“从构成 {[}本书{]}\ldots 之基础的诸原理导出，以致我不相信任何别的几何理论能以同时更直接、更简单的方式导出它们。如果考虑到图形的射影性质必然是图形所能具有的性质中最一般的一类，人们对这一点就丝毫不会感到惊讶；因而它们必定把一切其余的特殊性质或广延关系作为简单的推论包含在内”（（XII），第 248 页）。说实在的，在这番宣言之后，人们看到他处理角的问题时竟绕了很大的弯子——把角与二次曲线的焦点性质挂靠在一起，而不是直接让圆点发挥作用——不免有些诧异；而事实上，直到 30 年后，拉盖尔（当时还是综合理工学校的学生）才借助两条直线与同原点的两条迷向直线的交比，给出两直线间角的表达式（参见 §10，习题 5）（（XXI），卷 II，第 13 页）。最后，在凯莱 (XVIII d) 那里，一个基本思想被清楚地表达出来：平面图形的“度量”性质无非是该图形添上圆点之后的“射影”性质——这是通向“埃尔朗根纲领”的决定性里程碑。
\item
  双曲非欧几何诞生于 1830 年前后，起初多少置身于我们正概述其主线的那场运动之外。这一新几何源于本质上属于逻辑方面、触及经典几何基础的关切，其发明者们 (*) 以与欧几里得几何相同的公理化“综合”形式呈现它，且与射影几何毫无联系（在这种几何中，按照经典模式引入射影几何甚至看起来先验地就被排除了，因为唯一平行线的概念在这种几何中消失了）；大概正是由于这个原因，在很长时间里它很少引起法国、德国与英国的射影几何学派的兴趣。因此，当凯莱在上文引用的那篇基本著作 (XVIII d) 中想到把圆点（被视为一条“切向退化的”二次曲线）换成一条任意的二次曲线（他称之为“绝对形”）时，他完全没有想到把这个想法与罗巴切夫斯基–鲍耶几何联系起来，尽管他指出了他的构想如何给出两点间“距离”的新表达式，也尽管他提到了它与球面几何的联系。形势约在 1870 年前后发生变化，这时，继罗巴切夫斯基著作的传播，以及高斯著作与黎曼就职演讲的出版之后，非欧几何诸几何走到了数学现实的前沿。沿黎曼所开辟的道路，贝尔特拉米在不知道凯莱工作的情况下，于 1868 年重新发现了后者所给出的那些距离表达式，但所处的脉络完全不同：他把圆的内部视为一张常曲率曲面的像，其中测地线由直线表示 (XXIV)；两年之后，克莱因（独立于贝尔特拉米）作出了这些不同观点的综合，并以椭圆非欧空间的发现完成了这一综合 (XXV a) (**)。
\item
  在我们此处所考虑时期的后半段，出现了一个批判反思的时期，在此期间，“综合”几何的拥护者们不满足于已把坐标从证明中驱逐出去，还声称甚至在几何公理中也可以不用实数。这一学派的主要代表是冯·施陶特，他基本上完成了这一杰作 (XXIII)；此杰作在其时代乃至二十世纪相当长的时期内都备受推崇；尽管今天人们不再赋予这一层次的思想以同样的重要性——其富有成效的应用可能性已被证明相当贫乏——但仍须承认，冯·施陶特及其门徒们的努力，有助于澄清关于实或复“标量”在经典几何中所起作用的种种观念，并由此引入了在任意基体上的几何这一现代概念。
\end{enumerate}

约在 1860 年前后，“综合”几何正处于鼎盛时期，但其统治的终结已在迅速临近。在整个十八世纪始终笨重而不雅致的解析几何，在拉梅、博比利耶、柯西、普吕克与默比乌斯手中，终于获得了足以与对手平起平坐地竞争的优雅与简洁。尤其重要的是，约从 1850 年起，群与不变式的观念终于得到精确表述，逐渐占据了中心位置，并且人们认识到，经典几何的定理无非是相似变换群的不变式或协不变式之间恒等关系的表达 (*)，正如射影几何的定理表达射影群的协不变式之间的恒等式（即“合冲”）一样。这正是 F. 克莱因在其著名的“埃尔朗根纲领”(XXV b) 中以精湛技艺阐述的论点，他在其中主张放弃“综合”倾向与“解析”倾向之间徒劳无益的争论；他说，如果针对后者把特权角色给予任意坐标系的那一指责，“就以往使用坐标法时的缺陷方式而言往往确实是完全站得住脚的，那么一当涉及对这一方法的合理运用，它便不复成立\ldots{} 空间直观的领域并非解析法的禁地\ldots{}”；他还强调，“人们不可低估一个贴切的形式体系给后续研究带来的益处，因为它可以说替思考预先作了准备”（（XXV b），第 488–490 页）。

由此得到“几何学”诸定理按其源出之群的一种合理的“结构性”分类：射影几何对应线性群，度量问题对应正交群，“线性丛”的几何对应辛群。然而在这种无情的明晰之下，经典几何——代数几何与微分几何（**）除外，它们此时已确立为自治的科学——骤然枯萎，失去了全部光彩。基于使用变换的各种方法的一般化，早已使新定理的形成变得有些机械化：“今天，”夏尔于 1837 年在其 \emph{《历史概述》} 中说，“任何人都可以站出来，任取一条已知真理，把它置于各种一般变换原理之下；他将从中提取出别的真理，不同的或更一般的；而这些真理又可经受类似的运算；于是人们几乎可以无穷尽地增殖从第一条真理演绎出的新真理的数目\ldots{} 因此，在科学的现状下，谁愿意谁就可以在几何学中作出推广并有所创造；天才对于为大厦添砖加瓦已不再是不可或缺的了”（（XV b），第 268–269 页）。但随着不变量理论的进展，形势变得清楚得多：该理论终于（至少对“典型”群而言）成功地表述了一般方法，使人们原则上能以纯粹机械的方式写出一切代数协不变式及其全部“合冲”；这是一场胜利，它同时标志着经典不变量理论本身作为一个研究领域、以及“初等”几何的死亡——后者实际上已几乎沦为它的一本单纯词典。诚然，没有什么能使人们先验地预见，在可以如此随意涌现的无穷多条“定理”之中，哪些定理用适当的几何语言来表达会具有可与经典结果相比拟的简洁与优雅，而且仍存在一个有限的领域，在那里众多业余爱好者继续愉快地钻研（三角形几何、四面体几何、低次代数曲线与代数曲面的几何，等等）。但对职业数学家而言，矿藏已然枯竭，因为那里不再存在结构问题、不再存在能对数学其他部分产生反响的问题；群论与不变量理论的这一篇章，可以认为在另行通知之前业已关闭。(*)

\[
\ast \ast
\]

这样，在埃尔朗根纲领之后，欧氏几何与非欧几何从纯代数的观点来看，都成了表达双线性型理论结果的或繁或简的语言，而双线性型理论的进步则与不变量理论的进步携手并进 (*)。凡是涉及双线性型的秩的概念以及这些型与线性变换之间关系的一切内容，都由弗罗贝尼乌斯的工作最终澄清 (XXVII a)。自由 Z-模上交错型的典范表达式（§5，第 1 小节，定理 1）也应归功于弗罗贝尼乌斯 (XXVII b)；不过，斜对称行列式早在本世纪初就已出现在普法夫那里，与微分形式化为正规形相关；于 1827 年重新处理这一问题的雅可比 (XIX a)，知道奇数阶的斜对称行列式为零，并且正是他构造了普法夫式的表达式并证明它是偶数阶斜对称行列式的一个因子；但他未曾察觉后者是普法夫式的平方，这一点直到 1849 年才由凯莱建立 (XVIII b)。与二次型相伴的对称双线性型这一概念，是“极化”过程的最简单情形，而极化过程正是不变量理论的基本工具之一。以“标量积”之名，这一概念将享有巨大的成功，首先是在“向量分析”的普及者那里，然后从二十世纪起，得益于希尔伯特空间理论给它带来的意想不到的推广（见第 V 卷的历史注记）。也正是后一理论，将把算子的伴随这一概念显现出来（此前它几乎只是在线性微分方程理论中、以及在张量分析中通过协变与逆变指标在度量张量指挥棒下的舞蹈表现出来）；最后，也正是这一理论，将充分凸现埃尔米特型的概念，这一概念由埃尔米特于 1853 年在与数论研究相关的场合首先引入（（XXII），第 237 页），但直到 1925 年前后、直到复希尔伯特空间在量子理论中的应用，它一直多少处于数学主流的边缘。

另一方面，对正交群与相似变换群的研究——这两个群自十九世纪中叶起即被清楚地构想出来并被当作这样的对象处理，并且已成为二次型理论的核心——以及对其他“典型”群（线性群、辛群与酉群）的研究，扮演着越来越重要的角色。我们在此只能提一下这些群一方面在李群理论与微分几何中所起的关键作用，另一方面在二次型的算术理论中所起的关键作用（例如见 (XXXIII) 与 (XXXV)）(**)；正是由于这一情况，以及由于对偶概念向最形形色色问题的推广，才使得几乎没有哪一门现代数学理论不以这种或那种方式牵涉双线性型。无论如何我们必须指出，正是对（三维）旋转群的研究导致哈密顿发现了四元数 (XVII)；这一发现由 W. 克利福德加以推广，他于 1876 年引入了以他的名字命名的那些代数，并证明它们是四元数代数、或四元数代数与一个二次扩张的张量积 (XXVIII)。这些代数四年后被李普希茨重新发现 (XXIX)，他用它们给出了 n 个变量的正交变换的参数表示（推广了凯莱借助四元数理论对 n = 3 (XVIII a) 与 n = 4 (XVIII c) 所得到的那些表示（参见 §9，习题 15 与 16）），这些代数以及由之导出的“旋量”概念（见 (XXXII) 与 (XXXIV)），也因其在量子理论中的应用而在现代时期享有巨大的声誉。

最后，关于导致在半双线性型理论中几乎完全放弃对标量环任何限制的思想演变，我们只剩一句话要说——这是全部现代代数所共有的倾向，但在这里，它也许比在别处显现得更早。我们已经指出过在复数域上引入几何的丰硕成果（此外，在整个十九世纪，这种几何与实几何之间的混淆从未间断，有时甚至是危险的）；这里的明晰首先来自十九世纪末关于几何基础的那些公理化研究 (XXX)。在这些考察的过程中，希尔伯特及其后继者们特别是在研究各条公理之间的关系时，被引导去构造适当的反例，其中“基体”（交换或非交换）具有或多或少病态的性质，他们由此使数学家们习惯了全新类型的“几何”。从解析的观点看，伽罗瓦已经考虑过系数与变量都取值于一个有限素域的线性变换（（XVI），第 27 页）；在发展这些思想时，若尔当 (XXVI) 很自然地被引导去考虑这些体上的典型群，这些群的介入显现在数学的各个不同领域。迪克森约在 1900 年把若尔当的研究推广到一切有限域，而最近人们认识到，若尔当–迪克森理论的一大部分可以推广到绝对任意的“基体”的情形；这本质上应归功于迷向向量的一般性质与维特定理，它们在经典情形是平凡的，但直到 1936 年才对任意基体建立起来 (XXXI) (*)。

但是，在把半双线性型的研究推向越来越高的“抽象”时，把在 2 维与 3 维空间情形下来自经典几何的术语原封不动地保留下来，并把它推广到 n 维情形乃至无限维空间，被证明是极富启示性的。作为一门自主而有生命的科学已然过时的经典几何，由此被点化为当代数学的一种通用语言，其灵活与便利无可比拟。

